\subsection{Weyl inequalities}

In this issue, we consider Hilbert spaces over K = C.

Let E be a Hilbert space. For every \(n \in N\) , we recall that the Hilbertian exterior power \(\hat{\Lambda}^{n}E\) was defined in EVT, V, p. 34. For every Hilbert space F and for \(u \in \mathcal{L}(E;F)\) , we have also defined the linear map \(\hat{\Lambda}^{n}u \in \mathcal{L}(\hat{\Lambda}^{n}E;\hat{\Lambda}^{n}F)\) (loc. cit.). These constructions are functorial: for every Hilbert space G and for every linear map \(v \in \mathcal{L}(F;G)\) , the formula

\[
\widehat {\Lambda} ^ {n} v \circ \widehat {\Lambda} ^ {n} u = \widehat {\Lambda} ^ {n} (v \circ u)
\]

is valid (loc. cit., formula (28)).

Let H be a closed subspace of E and let \(i_{H}\) be the canonical injection from H into E. For every endomorphism u of E, we denote by \(u_{H}\) the endomorphism \(i_{H}^{*}ui_{H}\) of H.

Lemma 2. --- Let \(n \in N\) . The map \(\hat{\Lambda}^{n}i_{H}\) is a linear isometric mapping from \(\hat{\Lambda}^{n}H\) into \(\hat{\Lambda}^{n}E\) .

Let \((e_{j})_{j\in\mathrm{J}}\) be an orthonormal basis of H and \((e_{j})_{j\in\mathrm{J}^{\prime}}\) an orthonormal basis of E, where \(J\subset J^{\prime}\) . Equip \(J^{\prime}\) with a total order. The elements \(e_{j_{1}}\wedge\cdots\wedge e_{j_{n}}\) for \(j_{1}<\cdots<j_{n}\) in \(J^{\prime}\) (resp. in J) form an orthonormal basis of \(\widehat{\Lambda}^{n}E\) (resp. of \(\widehat{\Lambda}^{n}H\) ) by prop. 5 of EVT, V, p. 34, prop. 5. The lemma follows.

In what follows, we shall identify \(\hat{\Lambda}^n\mathrm{H}\) with a closed subspace of \(\hat{\Lambda}^n\mathrm{E}\) by means of the mapping \(\hat{\Lambda}^n i_{\mathrm{H}}\) .

Lemma 3. --- Let F be a Hilbert space and \(u \in \mathcal{L}(\mathrm{E};\mathrm{F})\). Let \(n \in \mathbf{N}\) . a) We have ( \(\widehat{\Lambda}^n u\) ) \(^*\) = \(\widehat{\Lambda}^n (u^*)\) ;

\begin{enumerate}
\def\labelenumi{\alph{enumi})}
\setcounter{enumi}{1}
\item
  If F = E and u is Hermitian (resp. positive, normal, unitary), then \(\hat{\Lambda}^{n}u\) is Hermitian (resp. positive, normal, unitary);
\item
  We have \(|\widehat{\Lambda}^n u| = \widehat{\Lambda}^n |u|\) ;
\item
  Let H be a closed subspace of E. The restriction \((\widehat{\Lambda}^{n}u)|\widehat{\Lambda}^{n}\mathrm{H}\) of \(\widehat{\Lambda}^{n}u\) to \(\widehat{\Lambda}^{n}\mathrm{H}\) is equal to \(\widehat{\Lambda}^{n}(u|\mathrm{H})\) (equality in \(\mathcal{L}(\widehat{\Lambda}^{n}\mathrm{H};\widehat{\Lambda}^{n}\mathrm{F})\) );
\item
  Suppose F = E. Let H be a closed subspace of E. We have \((\hat{\Lambda}^{n}u)_{\hat{\Lambda}^{n}\mathrm{H}}=\hat{\Lambda}^{n}(u_{\mathrm{H}})\) in \(\mathcal{L}(\hat{\Lambda}^{n}\mathrm{H})\) .
\end{enumerate}

The assertion \(a\) ) follows from EVT, V, p. 39, formula (13). It immediately implies that \(\widehat{\Lambda}^{n}u\) is Hermitian (resp. unitary) when \(u\) is Hermitian (resp. unitary). If \(u\) is positive, then \(u^{1/2}\) is Hermitian, hence so is \(\widehat{\Lambda}^{n}(u^{1/2})\) ; the relation \(\widehat{\Lambda}^{n}u = (\widehat{\Lambda}^{n}u^{1/2})^{2}\) implies that \(\widehat{\Lambda}^{n}u\) is positive (I, p. 138, prop. 8). This proves \(b\) ).

What precedes allows one to compute that

\[
\begin{array}{r l} & {\left(\widehat {\Lambda} ^ {n} | u |\right) ^ {2} = \widehat {\Lambda} ^ {n} (| u | ^ {2}) = \widehat {\Lambda} ^ {n} (u ^ {*} u)} \\ & {\qquad = \widehat {\Lambda} ^ {n} u ^ {*} \widehat {\Lambda} ^ {n} u = \left(\widehat {\Lambda} ^ {n} u\right) ^ {*} (\widehat {\Lambda} ^ {n} u) = | \widehat {\Lambda} ^ {n} u | ^ {2}.} \end{array}
\]

Since \(\widehat{\Lambda}^n |u|\) is positive by \(b\) ), it follows from prop. 16 of I, p. 118 that \(\widehat{\Lambda}^n |u| = |\widehat{\Lambda}^n u|\) , whence \(c\) ).

Let \(i_{\mathrm{H}}\) be the canonical injection of H into E. The space \(\widehat{\Lambda}^{n}\mathrm{H}\) is identified with a closed subspace of \(\widehat{\Lambda}^{n}\mathrm{E}\) by the mapping \(\widehat{\Lambda}^{n}i_{\mathrm{H}}\) , whence

\[
(\widehat {\Lambda} ^ {n} u) | \widehat {\Lambda} ^ {n} \mathrm{H} = \widehat {\Lambda} ^ {n} u \circ \widehat {\Lambda} ^ {n} i _ {\mathrm{H}} = \widehat {\Lambda} ^ {n} (u \circ i _ {\mathrm{H}}) = \widehat {\Lambda} ^ {n} (u | \mathrm{H}),
\]

and

\[
(\widehat {\Lambda} ^ {n} u) _ {\widehat {\Lambda} ^ {n} \mathrm{H}} = (\widehat {\Lambda} ^ {n} i _ {\mathrm{H}}) ^ {*} \circ \widehat {\Lambda} ^ {n} u \circ \widehat {\Lambda} ^ {n} i _ {\mathrm{H}} = \widehat {\Lambda} ^ {n} (i _ {\mathrm{H}} ^ {*} u i _ {\mathrm{H}}) = \widehat {\Lambda} ^ {n} (u _ {\mathrm{H}}),
\]

which proves \(d\) ) and \(e\) ).

Let F be a Hilbert space, and let \(u \in \mathcal{L}^{c}(\mathrm{E};\mathrm{F})\) . As in number 3 of IV, p. 151, we denote by \(I_{E}\) the set of dimensions of finite-dimensional subspaces F of E such that \(F \neq E\) . We recall that \((\alpha_{n}(u))_{n \in I_{E}}\) denotes the extended sequence of singular values of u (def. 3 of IV, p. 156).

PROPOSITION 10. --- One has the equality

\[
\prod_ {i = 0} ^ {n} \alpha_ {i} (u) = \left\| \widehat {\Lambda} ^ {n + 1} u \right\|
\]

for all \(n \in I_{E}\) .

Lemma 3, c) shows that \(\|\widehat{\Lambda}^{n+1}u\|=\|\widehat{\Lambda}^{n+1}|u\|\) ; moreover, since \(\alpha_{i}(u)=\alpha_{i}(|u|)\) for all \(i\in I_{E}\) (remark 5 of IV, p. 157, (2)), it suffices to prove the assertion for \(|u|\) . This allows us to assume that F=E and that u is positive.

Let then \(\mathrm{B}=(e_{j})_{j\in\mathrm{J}}\) be an orthonormal basis of E such that u is diagonalizable in the basis B (theorem 1 of IV, p. 149), and let \((\lambda_{j})_{j\in\mathrm{J}}\) be the family of eigenvalues of u in the basis B. Endow J with a total order, and denote by \(J_{n}\) the set of strictly increasing families \((j_{0},\ldots,j_{n})\in\mathbf{J}^{n+1}\) of elements of J. The vectors

\[
e _ {\iota} = e _ {j _ {0}} \wedge \dots \wedge e _ {j _ {n}}
\]

for \(\iota=(j_{0},\ldots,j_{n})\in\mathrm{J}_{n}\) form an orthonormal basis \(B_{n}\) of \(\widehat{\Lambda}^{n+1}\mathrm{E}\) (EVT, V, p. 34, prop. 5). For every \(\iota=(j_{0},\ldots,j_{n})\in\mathrm{J}_{n}\) , denote

\[
\lambda_ {\iota} = \prod_ {j = 0} ^ {n} \lambda_ {i _ {j}}.
\]

It follows that \((\hat{\Lambda}^{n + 1}u)e_{\iota} = \lambda_{\iota}e_{\iota}\) , hence \(\hat{\Lambda}^{n + 1}u\) is diagonal in the basis \(\mathrm{B}_n\) . Consequently, one has

\[
\| \widehat {\Lambda} ^ {n + 1} u \| = \sup _ {\iota \in \mathrm{I} _ {n}} \lambda_ {\iota}
\]

(lemma 1 of IV, p. 147), which is equal to the product \(\lambda_{0}(u)\cdots\lambda_{n}(u)\) of the \(n+1\) largest eigenvalues of \(u\) . The desired formula follows since \(\lambda_{i}(u)=\alpha_{i}(u)\) for all \(i\in\mathrm{I}_{\mathrm{E}}\) when \(u\) is positive (remark 5 of IV, p. 157, (3)).

In particular, if \(\alpha_{1}(u) < \alpha_{0}(u) = \| u\|\) , one sees that the inequality \(\| \widehat{\Lambda}^{n}(u)\| \leqslant \| u\|^{n}\) (EVT, V, p. 34, formula (29)) is not an equality in general if \(n \geqslant 2\) .

COROLLARY. --- Let G be a Hilbert space and \(v \in \mathcal{L}^{\mathrm{c}}(\mathrm{F};\mathrm{G})\) . We have

\[
\prod_ {i = 0} ^ {n} \alpha_ {i} (v \circ u) \leqslant \prod_ {i = 0} ^ {n} \alpha_ {i} (u) \alpha_ {i} (v)
\]

for all \(n \in \mathrm{I}_{\mathrm{E}}\) .

It suffices to note that

\[
\| \widehat {\Lambda} ^ {n + 1} (v u) \| = \| \widehat {\Lambda} ^ {n + 1} v \circ \widehat {\Lambda} ^ {n + 1} u \| \leqslant \| \widehat {\Lambda} ^ {n + 1} v \| \| \widehat {\Lambda} ^ {n + 1} u \|
\]

and to apply proposition 10.

Lemma 4. --- Let A be a ring. Let \(n \in \mathbf{N}\) and \((a_i)_{0 \leqslant i \leqslant n}\) and \((b_i)_{0 \leqslant i \leqslant n}\) be families of elements of A. For \(0 \leqslant j \leqslant n\) set

\[
\mathrm{A} _ {j} = \sum_ {i = 0} ^ {j} a _ {i}.
\]

One has

\[
\sum_ {i = 0} ^ {n} a _ {i} b _ {i} = \mathrm{A} _ {n} b _ {n} - \sum_ {i = 0} ^ {n - 1} \mathrm{A} _ {i} (b _ {i + 1} - b _ {i}).
\]

Set \(A_{-1}=0\) . We have \(a_{i}=A_{i}-A_{i-1}\) for \(0\leqslant i\leqslant n\) , hence

\[
\sum_ {i = 0} ^ {n} a _ {i} b _ {i} = \sum_ {i = 0} ^ {n} (\mathrm{A} _ {i} - \mathrm{A} _ {i - 1}) b _ {i} = \sum_ {i = 0} ^ {n - 1} \mathrm{A} _ {i} (b _ {i} - b _ {i + 1}) + \mathrm{A} _ {n} b _ {n},
\]

as desired.

Let I be an interval of \(\overline{R}\) not containing \(+\infty\) . Recall (FVR, I, p. 38, remark) that a continuous function \(f: I \to [-\infty, +\infty[\) is said to be convex if its restriction to the interior of I is a convex function with values in R; it then has a limit (finite or infinite) on the right at inf I. In the statements below, the product of 0 and an element of \(\{-\infty, +\infty\}\) is by convention equal to 0.

Lemma 5. --- Let I ⊂ R be an interval not containing +∞ and f an increasing convex function from I into \([-\infty, +\infty[\).

Let n be a natural integer and

\[
a _ {0} \geqslant a _ {1} \geqslant \dots \geqslant a _ {n}, \quad b _ {0} \geqslant b _ {1} \geqslant \dots \geqslant b _ {n}
\]

elements of I.

Let \(\varrho_{0} \geqslant \varrho_{1} \geqslant \cdots \geqslant \varrho_{n}\) be positive real numbers. Suppose that

\[
\sum_ {i = 0} ^ {j} a _ {i} \leqslant \sum_ {i = 0} ^ {j} b _ {i}.\tag{5}
\]

for every integer j such that \(0 \leqslant j \leqslant n\) . We then have

\[
\sum_ {i = 0} ^ {n} \varrho_ {i} f (a _ {i}) \leqslant \sum_ {i = 0} ^ {n} \varrho_ {i} f (b _ {i}).\tag{6}
\]

Suppose first that \(a_i \in \mathring{\mathrm{I}}\) , and hence \(f(a_i) \in \mathbf{R}\) , for all \(i\). Let \(i\) be such that \(0 \leqslant i \leqslant n\) , and let \((\alpha_i, \beta_i)\) be real numbers such that the line with equation \(y = \alpha_i x + \beta_i\) is a supporting line of the graph of \(f\) at the point \((a_i, f(a_i))\) (FVR, I, p. 37). One consequently has \(f(a_i) = \alpha_i a_i + \beta_i\) and \(f(b_i) \geqslant \alpha_i b_i + \beta_i\) since the graph of \(f\) lies above the supporting line, whence

\[
f (b _ {i}) - f (a _ {i}) \geqslant \alpha_ {i} (b _ {i} - a _ {i}).\tag{7}
\]

Moreover, if \(i < n\) , we have

\[
\alpha_ {i} \geqslant f _ {g} ^ {\prime} (a _ {i}) \geqslant f _ {d} ^ {\prime} (a _ {i + 1}) \geqslant \alpha_ {i + 1}
\]

(loc. cit. and FVR, I, p. 36, cor. 1); furthermore \(\alpha_{i} \geqslant 0\) since \(f\) is increasing (FVR, I, p. 22, corollary), whence \(\varrho_{i}\alpha_{i} \geqslant \varrho_{i+1}\alpha_{i+1} \geqslant 0\) for \(0 \leqslant i < n\) .

Let \(j\) be an integer such that \(0 \leqslant j \leqslant n\) . Set

\[
\mathrm{A} _ {j} = \sum_ {i = 0} ^ {j} (b _ {i} - a _ {i}),
\]

so that \(A_{j} \geqslant 0\) by hypothesis (5). Applying inequality (7) and then Lemma 4, we deduce

\[
\begin{array}{l} \sum_ {i = 0} ^ {n} \varrho_ {i} (f (b _ {i}) - f (a _ {i})) \geqslant \sum_ {i = 0} ^ {n} \varrho_ {i} \alpha_ {i} (b _ {i} - a _ {i}) \\ \qquad = \varrho_ {n} \alpha_ {n} A _ {n} + \sum_ {j = 0} ^ {n - 1} (\varrho_ {j} \alpha_ {j} - \varrho_ {j + 1} \alpha_ {j + 1}) A _ {j} \geqslant 0. \end{array}
\]

Consider the general case, and reason by induction on \(n\). If one of the \(a_{i}\) does not belong to the interior of I, we necessarily have \(a_{0} = \sup I\) or \(a_{n} = \inf I\) .

Suppose first that \(a_{n}=\inf I\) . We then have \(a_{n}\leqslant b_{n}\) and therefore \(\varrho_{n}f(a_{n})\leqslant\varrho_{n}f(b_{n})\) . The induction hypothesis, applied to the families \((a_{0},\ldots,a_{n-1})\) , \((b_{0},\ldots,b_{n-1})\) and \((\varrho_{0},\ldots,\varrho_{n-1})\) , implies

\[
\sum_ {i = 0} ^ {n - 1} \varrho_ {i} f (a _ {i}) \leqslant \sum_ {i = 0} ^ {n - 1} \varrho_ {i} f (b _ {i}),
\]

whence the desired inequality follows by adding \(\varrho_{n}f(a_{n})\). The case where \(a_{0} = \sup I\) is treated in a similar manner.

PROPOSITION 11 (Weyl inequalities). --- Let G be a Hilbert space and \(v \in \mathcal{L}^{\mathrm{c}}(\mathrm{F};\mathrm{G})\). Let \(g: \mathbf{R}_{+} \to [-\infty, +\infty[\) be a nondecreasing function such that the function \(g \circ \exp\) is convex. We have

\[
\sum_ {i = 0} ^ {n} g \left(\alpha_ {i} (v \circ u)\right) \leqslant \sum_ {i = 0} ^ {n} g \left(\alpha_ {i} (v) \alpha_ {i} (u)\right)
\]

for all \(n \in I_{E} \cap I_{F}\) .

Set \(I = [-\infty, +\infty[\) and \(f = g \circ \exp\). It is permissible to apply Lemma 5 with

\[
a _ {i} = \log (\alpha_ {i} (v \circ u)) \in \mathrm{I}, \quad b _ {i} = \log (\alpha_ {i} (v) \alpha_ {i} (u)) \in \mathrm{I}
\]

and \(\varrho_{i} = 1\) for \(0\leqslant i\leqslant n\) since

\[
\sum_ {i = 0} ^ {j} a _ {i} = \log \Bigl (\prod_ {i = 0} ^ {j} \alpha_ {i} (v \circ u) \Bigr) \leqslant \log \Bigl (\prod_ {i = 0} ^ {j} \alpha_ {i} (v) \alpha_ {i} (u) \Bigr) = \sum_ {i = 0} ^ {j} b _ {i}
\]

for \(0 \leqslant j \leqslant n\) by the corollary of prop. 10. Inequality (6) is then the desired conclusion.

Lemma 6. --- Let g and h be convex functions defined on intervals I and J of R, respectively. If g is increasing and defined on the image of h, then the function g ∘ h is convex on J.

Indeed, for \(t \in [0,1]\) and \((x,y) \in \mathrm{J} \times \mathrm{J}\) , we have

\[
\begin{array}{c} g (h (t x + (1 - t) y)) \leqslant g (t h (x) + (1 - t) h (y)) \\ \leqslant t g (h (x)) + (1 - t) g (h (y)). \end{array}
\]

COROLLARY. --- Let G be a Hilbert space and \(v \in \mathcal{L}^{\mathrm{c}}(\mathrm{F};\mathrm{G})\). Let \(n \in I_{E} \cap I_{F}\) .

\begin{enumerate}
\def\labelenumi{\alph{enumi})}
\tightlist
\item
  Let \(r \in \mathbf{R}_{+}^{*}\) . We have
\end{enumerate}

\[
\sum_ {i = 0} ^ {n} \alpha_ {i} (v \circ u) ^ {r} \leqslant \sum_ {i = 0} ^ {n} \alpha_ {i} (v) ^ {r} \alpha_ {i} (u) ^ {r}.
\]

\begin{enumerate}
\def\labelenumi{\alph{enumi})}
\setcounter{enumi}{1}
\tightlist
\item
  Let \(p, q, r \in \mathbf{R}_{+}^{*}\) such that \(\frac{1}{p} + \frac{1}{q} = \frac{1}{r}\) . Then
\end{enumerate}

\[
\left(\sum_ {i = 0} ^ {n} \alpha_ {i} (v \circ u) ^ {r}\right) ^ {1 / r} \leqslant \left(\sum_ {i = 0} ^ {n} \alpha_ {i} (v) ^ {p}\right) ^ {1 / p} \left(\sum_ {i = 0} ^ {n} \alpha_ {i} (u) ^ {q}\right) ^ {1 / q}.
\]

\begin{enumerate}
\def\labelenumi{\alph{enumi})}
\setcounter{enumi}{2}
\tightlist
\item
  Suppose that F = E. For every integer \(m \geqslant 2\) , we have
\end{enumerate}

\[
\sum_ {i = 0} ^ {n} \alpha_ {i} (u ^ {m}) \leqslant \left(\sum_ {i = 0} ^ {n} \alpha_ {i} (u) ^ {2}\right) ^ {m / 2}.
\]

Let \(r \in R_{+}^{*}\) and let g be the function defined on \(\mathbf{R}_{+}\) by \(g(x) = x^{r}\) . The function \(g \circ \exp\) is convex (lemma 6), so assertion a) follows from prop. 11 applied to the function g.

The assertion \(b\) ) follows from \(a\) ) and from Hölder's inequality (INT, I, prop. 4).

Let \(m \geqslant 2\) be an integer. Let us apply \(b)\) with \(r = 1\) , \(p = q = 2\) and \(v = u^{m-1}\). We find

\[
\sum_ {i = 0} ^ {n} \alpha_ {i} (u ^ {m}) \leqslant \left(\sum_ {i = 0} ^ {n} \alpha_ {i} (u ^ {m - 1}) ^ {2}\right) ^ {1 / 2} \left(\sum_ {i = 0} ^ {n} \alpha_ {i} (u) ^ {2}\right) ^ {1 / 2}.
\]

Let us prove c) by induction on \(m \geqslant 2\) . The preceding inequality establishes the assertion when m = 2. Suppose that \(m \geqslant 3\) and that the assertion is valid for m - 1; since we have

\[
\sum_ {i = 0} ^ {n} \alpha_ {i} (u ^ {m - 1}) ^ {2} \leqslant \left(\sum_ {i = 0} ^ {m} \alpha_ {i} (u ^ {m - 1})\right) ^ {2},
\]

the above inequality implies

\[
\sum_ {i = 0} ^ {n} \alpha_ {i} (u ^ {m}) \leqslant \left(\sum_ {i = 0} ^ {n} \alpha_ {i} (u ^ {m - 1})\right) ^ {1 / 2} \left(\sum_ {i = 0} ^ {n} \alpha_ {i} (u) ^ {2}\right) ^ {1 / 2} \leqslant \left(\sum_ {i = 0} ^ {n} \alpha_ {i} (u) ^ {2}\right) ^ {m / 2}
\]

by applying the induction hypothesis.

